\section*{How to Read This Paper}
\addcontentsline{toc}{section}{How to Read This Paper}

This paper began with two circuit observations that looked unrelated.

The first was constructive: in a coherent version of quantum teleportation,
two qubits end in a known state and can be cleaned for reuse. The second was
adversarial: in an educational spy-detector circuit, a more privileged
modification could restore the computational-basis output pattern while leaving
hidden correlations elsewhere in the state.

The connection is simple once it is stated. A qubit that becomes available is
both a resource and a boundary. A compiler may reuse it because a postcondition
says the earlier computation is finished with it. A malicious or untrusted
system layer may want that same qubit because it is available workspace. The
paper asks what must be proved before reuse is safe, and what a user must measure
before saying that a modified computation is actually equivalent to the one they
intended to run.

\paragraph{How to navigate the argument.}
Readers interested primarily in compiler optimization can begin with
\cref{sec:preliminaries,sec:case-study-i}. Readers interested in circuit-level
security can begin with \cref{sec:case-study-ii,sec:equivalence-levels}.
\Cref{sec:information-boundary} gives the general mathematical boundary, while
\cref{sec:threat-model,sec:capability-matrix,sec:defenses,sec:research-agenda}
develop the systems implications. The appendices collect resource counts, reproducibility details, a hyperlinked
cell-by-cell capability registry, the software artifact, and the questions that
remain open before submission.

\paragraph{Evidence labels.}
The paper separates what is proved from what is reconstructed or proposed:

\begin{center}
\begin{tabular}{@{}ll@{\qquad}ll@{}}
\evidencebadge{BaselineBlue}{PROVED} & symbolic derivation or theorem &
\evidencebadge{AdvancedTeal}{COMPUTED} & reproducible numerical result \\
\evidencebadge{InterceptOrange}{RECONSTRUCTED} & inferred from supplied exports &
\evidencebadge{ProposalPurple}{PROPOSED} & research hypothesis or design \\
\end{tabular}
\end{center}

\paragraph{Status and scope.}
This is Research Draft 0.7. It combines symbolic derivations, supplied Quirk
column arrays, complete five-wire amplitude exports, fixed-input branch analysis, exact reduced-state
calculations, ideal simulations, and a continuously tested software artifact. It does not claim a new teleportation
protocol, a break of BB84, a demonstrated attack on production quantum hardware,
or the recovery of arbitrary encrypted information.

The central claim is narrower: a known-state qubit should not be considered safe
to reuse merely because one visible output distribution looks correct. Safe
reuse requires a statement about state, separability, liveness, and information
flow.

\paragraph{Plain-language summary.}
A cleaned qubit is like an erased whiteboard. It is useful because the next
calculation can write on it. But ``the board looks blank'' and ``nothing secret
remains connected to the board'' are not the same statement.

Quantum mechanics makes the distinction unusually sharp. Classical-looking,
orthogonal information can be copied into hidden workspace while the original
register remains unchanged. An arbitrary unknown superposition cannot be
learned perfectly in the same way without leaving a physical trace. The hard
engineering problem lives between those extremes: real programs mix quantum
coherence with classical controls, compilers transform circuits across layers,
and users often verify only a small subset of the final state.
